% ============================================================
% 30_layout_spacing.tex — Layout-Grundlagen, Spacing-Skala, Box-Bindungen
% ============================================================

\setlength{\parindent}{0pt}

\widowpenalty=10000
\clubpenalty=10000
\brokenpenalty=4000
\predisplaypenalty=10000

\raggedcolumns
\newcounter{zsfFirstChap}\setcounter{zsfFirstChap}{1}

% ------------------------------------------------------------
% SPACING-SKALA (4 Stufen) — die einzige Quelle für Abstände
% ------------------------------------------------------------
% ── Dichte: die globale Stellschraube pro ZSF ────────────────────────
% Standard 1 = Referenzdichte des Templates. Für eine dichtere ZSF in
% preamble.tex VOR dem \input dieser Datei setzen, z.B.
%   \def\ZSFDensityFactor{0.85}
% Skaliert die vier Basisregister; alles, was aus ihnen berechnet wird, folgt
% automatisch (Box-Paddings, Balken-Polsterung, Section-Gaps, \columnsep,
% \tabcolsep, Display-Skips).
% Ersetzt das lokale Überschreiben von \ZSFspace* in einzelnen Kapiteln.
%
% BEWUSST AUSSERHALB (siehe rules/30_spacing → Reichweite):
%   \ZSFLeadingFactor  — Zeilenhöhe, eigener Regler (enthält die Schrift selbst)
%   \columnsep         — der Steg soll beim Verdichten nicht schrumpfen
%   \ZSFtitleBarHeight — Mindesthöhe des Titelstreifens, an die Schrift gebunden
%   \ZSFfigMaxHeight   — Bildhöhe ist Inhalt, kein Abstand
% Alles andere Vertikale hängt am Faktor. Wer etwas Neues hinzufügt, leitet es
% ebenfalls aus den Basisregistern ab — nicht als roher pt-Wert.
\providecommand{\ZSFDensityFactor}{1}
\newlength{\ZSFspaceXS}\setlength{\ZSFspaceXS}{\ZSFDensityFactor\dimexpr 1pt\relax}
\newlength{\ZSFspaceS} \setlength{\ZSFspaceS}{\ZSFDensityFactor\dimexpr 4pt\relax}
\newlength{\ZSFspaceM} \setlength{\ZSFspaceM}{\ZSFDensityFactor\dimexpr 7pt\relax}
\newlength{\ZSFspaceL} \setlength{\ZSFspaceL}{\ZSFDensityFactor\dimexpr 10pt\relax}

% Outer (Text-Box-Bindung) und Sep (Trennlinien) — beide folgen der Dichte.
% Outer ist bewusst als \ZSFspaceM abgeleitet statt erneut ausgeschrieben:
% es IST der Standard-Aussenabstand, kein eigener Wert.
\newlength{\ZSFspaceOuter}\setlength{\ZSFspaceOuter}{\ZSFspaceM}
\newlength{\ZSFspaceSep}  \setlength{\ZSFspaceSep}{\ZSFDensityFactor\dimexpr 3pt\relax}

% Absatzabstand im Fliesstext — der grösste einzelne vertikale Posten einer
% ZSF und deshalb der wichtigste Teil der Dichte-Stellschraube. Muss nach den
% Basisregistern stehen, damit \ZSFDensityFactor bereits bekannt ist.
\setlength{\parskip}{\ZSFDensityFactor\dimexpr 3.5pt\relax
             plus \ZSFDensityFactor\dimexpr 1pt\relax
            minus \ZSFDensityFactor\dimexpr 0.5pt\relax}

% ── Zeilenhöhe: die zweite globale Stellschraube ─────────────────────
% Bewusst NICHT an \ZSFDensityFactor gekoppelt. Abstände sind Leerraum und
% vertragen jede Skalierung; Zeilenhöhe enthält die Schrift selbst — sie mit
% demselben Faktor zu stauchen, lässt Formeln und Akzente kollidieren.
% Wer zusätzlich zur Dichte enger setzen will, dreht hier separat:
%   \def\ZSFLeadingFactor{0.92}   % in preamble.tex, vor diesem \input
\providecommand{\ZSFLeadingFactor}{0.95}
\linespread{\ZSFLeadingFactor}

% ------------------------------------------------------------
% Spalten-Layout
% ------------------------------------------------------------
% \columnsep folgt der Dichte bewusst NICHT: Beim Verdichten sollen die
% Zeilen enger stehen, nicht die Spalten zusammenrücken — ein schmalerer
% Steg macht die ZSF schlechter lesbar, obwohl er kaum Platz spart.
\newlength{\ZSFcolumnSep}\setlength{\ZSFcolumnSep}{10pt}
\setlength{\columnsep}{\ZSFcolumnSep}
\setlength{\multicolsep}{\ZSFspaceS}
\setlength{\emergencystretch}{2em}

% ------------------------------------------------------------
% Schwellwerte und Titelbalken-Hilfsgrössen
% ------------------------------------------------------------
% Mindesthöhe, die vor einem Titelbalken frei sein muss — sonst bricht die
% Spalte davor um. Die Reserve muss den Balken UND den Anfang der folgenden
% Box tragen, sonst bleibt der Balken allein am Spaltenende stehen.
% Deshalb staffeln sich die Werte nach der Grösse der eröffneten Einheit:
% Kapitel > Abschnitt > Unterabschnitt.
\newlength{\ZSFbreakThresholdChapter}\setlength{\ZSFbreakThresholdChapter}{10\baselineskip}
\newlength{\ZSFbreakThreshold}\setlength{\ZSFbreakThreshold}{7\baselineskip}
\newlength{\ZSFbreakThresholdSubsubsection}\setlength{\ZSFbreakThresholdSubsubsection}{5\baselineskip}
% Feinjustierung des Abstands direkt nach einem Kapitelbalken. Steht auf 0pt und
% sieht dadurch nach totem Code aus — ist es nicht: Das zugehörige
% \ZSFPostChapterSpacing (40_colors_structure) setzt ein \vspace*, und ein
% \vspace* ist auch bei 0pt nicht verwerfbar. Es erzwingt den Übergang in den
% vertikalen Modus und verankert den Absatzanfang nach dem Balken. Entfernt man
% es, verschiebt sich der Satz messbar (geprüft: alle drei Seiten ändern sich).
% Der Wert darf verstellt werden, die Zeile nicht wegfallen.
\newlength{\ZSFpostChapterNudge}\setlength{\ZSFpostChapterNudge}{0pt}
\newlength{\ZSFtitleBarHeight}\setlength{\ZSFtitleBarHeight}{9.2pt}
\newlength{\ZSFsubsectionBarBeforeSkip}\setlength{\ZSFsubsectionBarBeforeSkip}{\ZSFspaceL}
\newlength{\ZSFsubsubsectionBarBeforeSkip}\setlength{\ZSFsubsubsectionBarBeforeSkip}{\ZSFspaceS}
\newlength{\ZSFsubsectionAfterChapterGap}\setlength{\ZSFsubsectionAfterChapterGap}{\ZSFspaceXS}
\newlength{\ZSFsubsectionAfterGap}\setlength{\ZSFsubsectionAfterGap}{\ZSFspaceXS}
\newlength{\ZSFchapterBarAfterSkip}\setlength{\ZSFchapterBarAfterSkip}{\ZSFDensityFactor\dimexpr 2pt\relax}
\newlength{\ZSFboxArc}\setlength{\ZSFboxArc}{\ZSFspaceXS}
\newlength{\ZSFboxRuleWidth}\setlength{\ZSFboxRuleWidth}{0.5pt}
\newlength{\ZSFgoalConditionRuleWidth}\setlength{\ZSFgoalConditionRuleWidth}{1.5pt}
\newlength{\ZSFgoalConditionBorderlineWidth}\setlength{\ZSFgoalConditionBorderlineWidth}{0.2pt}
\newcommand{\ZSFtableArrayStretch}{1.04}

% Maximale Höhe für Bilder in figboxen (verhindert überhohe Aspect-Ratios).
\newlength{\ZSFfigMaxHeight}\setlength{\ZSFfigMaxHeight}{2.6cm}

% Abstand der TikZ-Spannklammer (\drawbrace, Opt-in 11_math_advanced) von der
% Grundlinie des annotierten Terms. Ohne ihn läuft die Klammer als waagrechte
% Linie mitten durch die Formel, weil die Marker auf der Grundlinie sitzen.
\newlength{\ZSFbraceDrop}\setlength{\ZSFbraceDrop}{\ZSFspaceM}

% Tokens für statementbox / procedure
\newlength{\ZSFstatementBarWidth}\setlength{\ZSFstatementBarWidth}{2pt}
\newlength{\ZSFprocStepGap}      \setlength{\ZSFprocStepGap}{\ZSFspaceXS}
\newlength{\ZSFprocItemSep}      \setlength{\ZSFprocItemSep}{\ZSFspaceXS}
\newlength{\ZSFprocLeftMargin}   \setlength{\ZSFprocLeftMargin}{1.6em}
\newlength{\ZSFlistLeftMargin}   \setlength{\ZSFlistLeftMargin}{1.2em}
\newlength{\ZSFlistLabelSep}     \setlength{\ZSFlistLabelSep}{0.4em}

% Kleinstes vertikales Innenmass für Text an einer Kante: Zeilenluft in
% Tabellen (\extrarowheight in 70_document_settings) und die knappe Polsterung
% von goalbox/valuegrid.
\newlength{\ZSFtextMinPad}\setlength{\ZSFtextMinPad}{\ZSFspaceXS}

% ── Horizontaler Innenabstand: eine Stellschraube für alle Inhaltsboxen ──
% Vorher stand an jeder Box ein eigener Wert (0pt/XS/S/M). Wer die Ränder
% ändern wollte, musste jede Box einzeln anfassen — genau der Fall, in dem
% eine KI repariert statt anzupassen.
%
% \ZSFboxPadX        — Standard-Rand zwischen Boxkante und Inhalt.
% \ZSFboxPadXBar     — für Boxen mit linkem Farbbalken: der Inhalt muss den
%                      Balken zusätzlich freihalten, deshalb links mehr.
% \ZSFtableEdgePad   — für Boxen, deren Hintergrund bis an die Kante laufen
%                      soll (Zebra, Wertegitter). Dort trägt nicht die Box
%                      den Abstand, sondern die äusserste Tabellenzelle.
\newlength{\ZSFboxPadX}      \setlength{\ZSFboxPadX}{\ZSFspaceS}
\newlength{\ZSFboxPadXBar}   \setlength{\ZSFboxPadXBar}{\ZSFspaceM}
\newlength{\ZSFtableEdgePad} \setlength{\ZSFtableEdgePad}{\ZSFspaceM}

% Vertikal dasselbe: Innenabstand oben/unten und der Abstand NACH einer Box.
% \ZSFBoxBeforeSkip gab es schon, das Gegenstück fehlte — deshalb stand
% "after skip=\ZSFspaceS" fünfmal einzeln im Code.
\newlength{\ZSFboxPadY}      \setlength{\ZSFboxPadY}{\ZSFspaceS}
\newlength{\ZSFboxPadYTight} \setlength{\ZSFboxPadYTight}{\ZSFspaceXS}
\newlength{\ZSFBoxAfterSkip} \setlength{\ZSFBoxAfterSkip}{\ZSFspaceS}

% Balken haben ihren eigenen vertikalen Rhythmus — bewusst getrennt vom
% Box-Innenabstand, damit sich Balkenhoehe und Boxpolsterung unabhaengig
% einstellen lassen.
\newlength{\ZSFbarPadY}      \setlength{\ZSFbarPadY}{\ZSFspaceS}
\newlength{\ZSFbarPadYTight} \setlength{\ZSFbarPadYTight}{\ZSFspaceXS}

% ------------------------------------------------------------
% Code-Längen — nur für die optionalen Code-Module benötigt
% (styles/65_code_style.tex + 67_code_comments.tex). Schadlos
% wenn die Module in preamble.tex deaktiviert sind.
% MaxWidth:   max. Breite des rechts-stehenden Kommentars
% Threshold:  Schwelle für Auto-Overflow (Code+Kommentar > Threshold -> oben)
% Default-Register sind immutabel und dienen \ResetCodeCommentThreshold
% als Anker; \SetCodeCommentThreshold modifiziert nur die "lebenden" Register.
% ------------------------------------------------------------
\newlength{\ZSFcodeMiniSkip}  \setlength{\ZSFcodeMiniSkip}{1mm}
\newlength{\ZSFcodeFrameRule} \setlength{\ZSFcodeFrameRule}{0.4pt}
\newlength{\ZSFcodeXMargin}   \setlength{\ZSFcodeXMargin}{8pt}
\newlength{\ZSFcodeCommentMaxWidthDefault}
  \setlength{\ZSFcodeCommentMaxWidthDefault}{.25\columnwidth}
\newlength{\ZSFcodeCommentThresholdDefault}
  \setlength{\ZSFcodeCommentThresholdDefault}{\maxdimen}
\newlength{\ZSFcodeCommentMaxWidth} \setlength{\ZSFcodeCommentMaxWidth}{\ZSFcodeCommentMaxWidthDefault}
\newlength{\ZSFcodeCommentThreshold}\setlength{\ZSFcodeCommentThreshold}{\ZSFcodeCommentThresholdDefault}
\newlength{\ZSFcodeCommentGap}      \setlength{\ZSFcodeCommentGap}{\ZSFspaceM}

% Valuegrid, Index, Readability und dokumentweite Overlays
\newlength{\ZSFvaluegridRowSep}       \setlength{\ZSFvaluegridRowSep}{\ZSFspaceSep}
\newlength{\ZSFreadabilityEmergencyStretch}\setlength{\ZSFreadabilityEmergencyStretch}{2.5em}
% Flatterzone im Fliesstext: schmaler als der ragged2e-Standard (0pt plus 2em),
% damit die schmalen 4 Spalten gut gefüllt bleiben. Wirkt nur über
% \ZSFReadableOn — Tabellen-Spalten (L/Y) und Index behalten den Standard.
\newlength{\ZSFreadabilityRaggedRightskip}\setlength{\ZSFreadabilityRaggedRightskip}{0pt plus 1em}
\newlength{\ZSFindexItemIndent}       \setlength{\ZSFindexItemIndent}{1.2em}
\newlength{\ZSFindexSubitemIndent}    \setlength{\ZSFindexSubitemIndent}{2em}
\newlength{\ZSFindexSubitemOffset}    \setlength{\ZSFindexSubitemOffset}{1em}
\newlength{\ZSFindexSubsubitemIndent} \setlength{\ZSFindexSubsubitemIndent}{2.6em}
\newlength{\ZSFindexSubsubitemOffset} \setlength{\ZSFindexSubsubitemOffset}{1.8em}
\newlength{\ZSFindexParStretch}       \setlength{\ZSFindexParStretch}{0.3pt}
\newlength{\ZSFfooterBottomInset}     \setlength{\ZSFfooterBottomInset}{1mm}
\newlength{\ZSFfooterRightInset}      \setlength{\ZSFfooterRightInset}{2mm}
% Freiraum am Satzspiegel-Fuss für das Footer-Overlay (Release-Kennung und
% Seitenzahl). Das Overlay zeichnet ausserhalb des Textflusses und reserviert
% deshalb KEINE Höhe — ohne diesen Freiraum überdruckt es die letzte Zeile
% einer voll gefüllten Spalte. Der Wert wird in 70_document_settings aus der
% Footer-Schrift gemessen und an geometry übergeben; hier steht nur das
% Register, damit alle Masse an einem Ort liegen.
\newlength{\ZSFfooterClearance}
\newlength{\ZSFitemizeLeftMargin}     \setlength{\ZSFitemizeLeftMargin}{3mm}
\newlength{\ZSFenumerateLeftMargin}   \setlength{\ZSFenumerateLeftMargin}{4mm}
\newlength{\ZSFidentityMarkerInset}   \setlength{\ZSFidentityMarkerInset}{2mm}

% ============================================================
% TEXT-BOX-BINDUNGEN
% Koppeln Fliesstext direkt an Boxen — verhindert visuellen Bruch.
% ============================================================

% Robustes Skip-/Penalty-Aufrollen (vmode-/hmode-aware)
\newcommand{\ZSFRobustUnskip}{%
  \ifvmode
    \ifdim\lastskip>0pt\relax\vskip-\lastskip\fi
    \ifnum\lastpenalty=0\else\unpenalty\fi
    \ifdim\lastskip>0pt\relax\vskip-\lastskip\fi
  \else
    \ifdim\lastskip>0pt\relax\unskip\fi
    \ifnum\lastpenalty=0\else\unpenalty\fi
    \ifdim\lastskip>0pt\relax\unskip\fi
  \fi
}

% \textVorBox / \textNachBox binden einen Satz optisch an die Box davor bzw.
% dahinter. INNERHALB einer formulabox binden dieselben zwei Makros an die
% Formelzeile — vorher gab es dafür \textVorFormel/\textNachFormel, also vier
% Namen für zwei Rollen (rules/02_ai_mandate → Katalog-Ökonomie).
% Die Umschaltung hängt an \ifZSF@inFormulaBox, das die formulabox in ihrem
% 'before upper' setzt; das Flag ist damit auf die Boxgruppe beschränkt.
\makeatletter
% Hier deklariert, nicht in 60_boxes: Das Flag gehört zu diesen beiden Makros
% und muss vor ihnen existieren, egal in welcher Reihenfolge die Module
% wachsen. Gesetzt wird es von der formulabox.
\newif\ifZSF@inFormulaBox
% Text der zur Box DARUNTER gehört (Einleitung)
\newcommand{\textVorBox}[1]{%
  \ifZSF@inFormulaBox
    \par
    \begingroup
    \setlength{\parskip}{0pt}%
    \noindent\raggedright\textcolor{MutedTextColor}{\textbf{#1}}\par
    \endgroup
    \vspace{\ZSFspaceXS}%
    \centering
  \else
    \par
    \addvspace{\ZSFspaceOuter}%
    \begingroup
    \setlength{\parskip}{0pt}%
    {\small\color{TextVorBoxColor} \noindent #1\par}%
    \endgroup
    \nopagebreak
    \vspace{\ZSFspaceXS}%
    \vspace{-\ZSFspaceS}%
    \nopagebreak
  \fi
}

% Text der zur Box DARÜBER gehört (Ergänzung/Folgerung)
\newcommand{\textNachBox}[1]{%
  \ifZSF@inFormulaBox
    \par
    \vspace{\ZSFspaceXS}%
    \begingroup
    \setlength{\parskip}{0pt}%
    \noindent\raggedright\textcolor{MutedTextColor}{\small #1}\par
    \endgroup
    \centering
  \else
    \par
    \ZSFRobustUnskip
    \vspace{\ZSFspaceXS}%
    \begingroup
    \setlength{\parskip}{0pt}%
    \nopagebreak
    {\small\color{TextNachBoxColor} \noindent #1\par}%
    \endgroup
    \addvspace{\ZSFspaceOuter}%
  \fi
}
\makeatother

% ------------------------------------------------------------
% Semantische Gap-Helfer für Kapitel (sparsam einsetzen)
% ------------------------------------------------------------
% \ZSFgap[<Stufe>] — ein Makro mit Regler statt vier Namen für eine Skala
% (rules/02_ai_mandate → Katalog-Ökonomie). Stufe ist XS, S, M oder L;
% ohne Angabe gilt M.
\NewDocumentCommand{\ZSFgap}{O{M}}{\vspace{\csname ZSFspace#1\endcsname}}
% \ZSFSectionGap bleibt als eigener Name, weil er etwas anderes ausdrückt als
% eine Grösse: den Themenwechsel. Dass er aktuell derselbe Abstand wie
% \ZSFgap[L] ist, ist eine Einstellung dieser Zeile, keine Verdopplung.
\newcommand{\ZSFSectionGap}{\ZSFgap[L]}

% ------------------------------------------------------------
% Vertikaler Einschub INNERHALB einer Box
% ------------------------------------------------------------
% Formelnote, Trennlinie, Item-Heading und Folgerung setzen alle denselben
% Einschub: laufenden \parskip zurücknehmen, definierten Abstand setzen.
% Das Muster stand neunmal von Hand im Code — und jede Kopie hatte \parskip
% fest verdrahtet, wodurch der Absatzabstand nicht an der Dichte-Stellschraube
% hängen konnte. Eine Stelle statt neun; wer den Rhythmus ändert, ändert hier.
\newcommand{\ZSFInterlude}[1]{\par\vspace{-\parskip}\vspace{#1}}
% Trenner-Abstand um den Mittelpunkt im Index-Locator "6.2 · (S. 4)" (66_index).
% Bewusst als Makro statt \newlength: das em soll an der Verwendungsstelle in der
% Index-Schrift ausgewertet werden. Eine Länge würde es in der Präambel-Schrift
% einfrieren und den Abstand dadurch messbar verschieben.
\newcommand{\ZSFindexLocatorSep}{\hspace{0.2em}}
